\documentclass[12pt,oneside]{article}
\usepackage[T1]{fontenc}
\usepackage{amssymb, amsmath, amsfonts}
\usepackage[mathscr]{euscript}
\usepackage{theorem}
\usepackage[english]{babel}
\usepackage{bm}
\usepackage[bookmarks=true,pdfauthor={Tomi Setsu},pdftitle={TOS: Actor-Oriented Blockchain}]{hyperref}
\usepackage{fancyhdr}
\usepackage{caption}
\usepackage{booktabs}
\usepackage{listings}
\usepackage{xcolor}
%
\setlength{\headheight}{15.2pt}
\pagestyle{fancy}
\renewcommand{\headrulewidth}{0.5pt}
%
\fancyhf{}
\fancyfoot[C]{\thepage}
\def\markbothsame#1{\fancyhead[C]{#1}}
\def\mysection#1{\section{#1}\fancyhead[C]{\textsc{Chapter \textbf{\thesection.} #1}}}
\def\mysubsection#1{\subsection{#1}\fancyhead[C]{\small{\textsc{\textrm{\thesubsection.} #1}}}}
%
\lstdefinelanguage{Tola}{
  keywords={actor, abstract, function, external, internal, view, payable, pure,
    returns, require, revert, emit, event, modifier, constructor,
    if, else, for, while, do, repeat, return, import, is,
    mapping, struct, enum, type, constant, immutable, try, catch,
    unchecked, unsafe, cell, interface, library, bounce, tick, tock,
    receive_external, selfdestruct, setCode},
  keywordstyle=\color{blue}\bfseries,
  ndkeywords={uint8, uint16, uint32, uint64, uint128, uint256,
    int8, int16, int32, int64, int128, int256,
    bool, bytes, bytes32, string, ActorId, SendMode, SendOptions, ReserveMode},
  ndkeywordstyle=\color{teal}\bfseries,
  sensitive=true,
  comment=[l]{//},
  morecomment=[s]{/*}{*/},
  commentstyle=\color{gray}\itshape,
  stringstyle=\color{red},
  morestring=[b]",
}
\lstset{
  language=Tola,
  basicstyle=\small\ttfamily,
  breaklines=true,
  frame=single,
  xleftmargin=1em,
  numbers=none,
}
%
\hfuzz=0.8pt

\title{\vspace{-1cm}\textbf{TOS: An Actor-Oriented Blockchain Platform}\\[0.5em]
\large Parallel Execution, Native Standard Library, and Developer-First Design}
\author{Tomi Setsu}
\date{April 2026}

\begin{document}

\maketitle

\begin{abstract}
TOS is a high-performance blockchain platform designed from the ground
up around the Actor Model. Every smart contract is a collection of
fine-grained, independently-executing actors, giving TOS intra-contract
parallelism that existing blockchains cannot achieve. The platform
combines a Cell/BOC data layer---providing native Merkle proofs and
compact serialization---with a simple key-value state interface exposed
through a new Solidity-like language called Tola. TOS ships a native
C++ standard actor library---analogous to Ethereum's OpenZeppelin but
compiled into the node binary---enabling zero-deployment, zero-audit
standard primitives for tokens, NFTs, governance, and DeFi.
\end{abstract}

\clearpage
\tableofcontents

\clearpage

%=============================================================
\mysection{Introduction and Motivation}\label{sect:intro}
%=============================================================

\mysubsection{The State of Blockchains in 2026}

The blockchain landscape in 2026 is shaped by five converging forces:

\begin{enumerate}
\item \textbf{The AI agent economy.} Autonomous AI agents are
  becoming first-class economic actors. By early 2026, agents hold
  their own wallets, negotiate payments via standardized HTTP payment
  protocols~\cite{x402}, and execute thousands of micro-transactions
  per second without human intervention. Account abstraction and tool
  invocation standards (such as the Model Context Protocol) are
  maturing rapidly. This demands blockchains that can handle
  \emph{massive volumes of small, autonomous, machine-to-machine
  transactions} with low latency and predictable fees---requirements
  that rollup-fragmented architectures and serial-execution VMs cannot
  meet.

\item \textbf{Throughput vs.\ decentralization.} High-throughput
  chains achieve speed by concentrating hardware requirements, while
  rollup-centric roadmaps fragment liquidity and user experience
  across hundreds of Layer-2 networks. Recent industry efforts toward
  parallel EVM execution and Verkle Tree state representations are
  promising, but remain incremental patches to fundamentally serial
  architectures designed a decade ago.

\item \textbf{Parallel execution vs.\ developer simplicity.}
  Object-based and resource-based parallel execution models have
  emerged, but their programming paradigms require developers to learn
  entirely new languages and mental models. Account-locking approaches
  require manual specification of all state touched by a transaction,
  creating friction and frequent failures. The vast majority of smart
  contract developers remain trained on Solidity-style syntax.

\item \textbf{Privacy as a first-class requirement.} Regulators
  demand transparency while users demand confidentiality. The
  zero-knowledge proof market is projected to reach \$7.6B by 2033.
  Industry efforts have produced encrypted rollup transactions,
  enterprise privacy sandboxes, and confidential compute frameworks
  for smart contracts. Yet no Layer-1 platform offers privacy as a
  \emph{native, per-actor capability} rather than an afterthought
  bolted onto a transparent execution model.

\item \textbf{The quantum threat.} Recent estimates suggest a
  cryptographically-relevant quantum computer could break elliptic
  curve signatures---the scheme securing the vast majority of
  blockchains---in under nine minutes~\cite{QuantumThreat}. A
  significant fraction of on-chain value has exposed public keys and
  would be immediately vulnerable. NIST standardized post-quantum
  signature schemes (SLH-DSA/FIPS~205, ML-DSA/FIPS~204) in 2024, and
  the industry broadly targets 128-bit provable post-quantum security
  by decade's end. Yet no shipping blockchain has a \emph{modular,
  actor-level mechanism} for upgrading cryptographic primitives
  without a network-wide hard fork.
\end{enumerate}

TOS addresses all five forces through a single architectural
principle: the \textbf{Actor Model} applied at the smart contract
level, with natural extension points for AI agent workloads, per-actor
privacy, and post-quantum authentication.

\mysubsection{The Actor Model Insight}

Two axioms from first principles:

\begin{enumerate}
\item Two execution units can run in parallel \emph{if and only if}
  their mutable state is completely disjoint.
\item Blockchain determinism requires all validators to agree on
  execution order and results.
\end{enumerate}

The conventional account model treats each smart contract as a
monolithic state blob. A popular token contract with millions of
users must serialize \emph{every} transfer because they all touch the
same balance mapping. Sharding helps when different contracts are
busy, but cannot help when one contract is the bottleneck.

The Actor Model solves this by decomposing a contract's state into
\textbf{independent actors}, each owning an isolated state partition
with its own logical timeline. Actors within the same contract
communicate via asynchronous messages, and---crucially---can execute
\textbf{in parallel} because their state never overlaps.

Crucially, the actor boundary is also a natural \textbf{privacy
boundary} and a natural \textbf{cryptographic boundary}: each actor
can independently adopt encrypted state, zero-knowledge proofs, or
post-quantum signature schemes without affecting other actors or
requiring a global protocol upgrade.

\mysubsection{Design Principles}

TOS is governed by seven design principles:

\begin{enumerate}
\item \textbf{Actor as the minimum execution unit.} Not accounts, not
  transactions---actors. Each actor owns its state, processes one
  message at a time, and communicates only through async messages.

\item \textbf{Cell inside, KV outside.} Internally, TOS stores all
  state in an elegant Cell/Merkle data structure that provides native
  Merkle proofs and compact consensus encoding. Developers see only
  \texttt{state\_get}/\texttt{state\_put}---a key-value interface.
  Zero-cost abstraction: the compiler maps KV operations to Cell
  dictionary operations transparently.

\item \textbf{Native standard library.} The most-used contract
  patterns (tokens, NFTs, access control, governance, DeFi primitives)
  ship as native C++ built-in actors, compiled into the node binary.
  Users instantiate actors via message, not deployment. Audited once,
  shared by all.

\item \textbf{Familiar syntax.} The Tola language uses widely-adopted
  smart contract syntax conventions (\texttt{function}, \texttt{require},
  \texttt{mapping}, \texttt{modifier}, \texttt{event}) so the largest
  developer community can onboard with minimal learning. The single
  new concept: cross-actor calls are \texttt{.send()}, not synchronous
  function calls.

\item \textbf{AI-agent ready.} Actors are the natural unit for agent
  wallets: each agent is an actor with its own state, spending policy,
  and message interface. The async messaging model matches the
  request/response pattern of MCP and x402, enabling high-throughput
  autonomous agent economies without blocking or reentrancy risks.

\item \textbf{Privacy-extensible.} Because each actor's state is an
  isolated HashmapE, a confidential actor can encrypt its state entries
  and prove correct execution via zero-knowledge proofs without
  touching any other actor's data. Privacy is a per-actor choice, not
  a network-wide mode switch.

\item \textbf{Quantum-resilient by design.} Actor-level authentication
  means each actor can independently adopt post-quantum signature
  schemes (e.g., SPHINCS+/SLH-DSA). A built-in \texttt{PostQuantumAuth}
  actor module can replace ECDSA for any actor without a chain-wide
  hard fork. The upgrade is granular: critical actors (treasuries,
  bridges) can go post-quantum immediately while the rest of the
  network migrates at its own pace.
\end{enumerate}

%=============================================================
\mysection{Architecture Overview}\label{sect:arch}
%=============================================================

\mysubsection{The Actor as the Atomic Unit}

The minimum addressable unit in TOS is not an account but an
\textbf{actor}. An account is simply a deployment container that holds
shared code and groups related actors together:

\begin{verbatim}
Account (unique address)
 +-- code: Cell             <- shared by all actors in this account
 +-- balance: Coins         <- shared (account-level)
 +-- actors: HashmapE       <- actor_id -> actor state (Cell dict)
 |    +-- actor_0: HashmapE <- independent KV state
 |    +-- actor_1: HashmapE <- independent KV state
 |    +-- ...
 +-- actor_lt: HashmapE     <- actor_id -> independent logical time
\end{verbatim}

Each actor is identified by a three-part address:
\texttt{workchain:account\_id:actor\_id}. The \texttt{actor\_id} is
deterministically derived: $\texttt{actor\_id} = \text{sha256}(\texttt{account\_address} \| \texttt{discriminator})$.

This design stands in contrast to conventional account-based chains
where a single contract holds one monolithic state tree behind one
logical timeline, forcing all operations to serialize.

\mysubsection{Parallel Execution}

Different actors under the same account execute in parallel because
their state (HashmapE) is disjoint. The node's \texttt{td::actor}
scheduler naturally handles this: each on-chain actor maps to a
\texttt{td::actor::Actor} instance, and the scheduler's thread pool
dispatches them concurrently. Same-actor messages remain serialized
(\texttt{ActorLocker} guarantee).

\begin{center}
\begin{tabular}{lcc}
\toprule
& \textbf{Account Model} & \textbf{Actor Model (TOS)} \\
\midrule
Minimum unit & Account & Actor \\
Parallelism & $N$ accounts & $N \times M$ actors \\
Hot-contract throughput & 1$\times$ & $M\times$ (actors per contract) \\
Reentrancy & Possible & Impossible by design \\
\bottomrule
\end{tabular}
\end{center}

\mysubsection{Cell Inside, KV Outside}

TOS uses the Cell data structure as its internal storage layer.
Every actor's state is a \texttt{HashmapE} (Cell-based dictionary).
This yields three properties:

\begin{itemize}
\item \textbf{Native Merkle proofs.} The proof chain
  \texttt{block\_root} $\to$ \texttt{shard\_state} $\to$
  \texttt{account} $\to$ \texttt{actor} $\to$ \texttt{key=value} is
  entirely Cell-based, giving every actor state entry a compact,
  verifiable Merkle path.
\item \textbf{Uniform BOC encoding.} Blocks, transactions, and
  messages all use Cell encoding, providing a single serialization
  format across the entire stack.
\item \textbf{Developers see KV.} The new TVM opcodes
  \texttt{STATEGET}/\texttt{STATESET}/\texttt{STATEDEL} wrap
  \texttt{HashmapE::get}/\texttt{set}/\texttt{remove}. Developers
  write \texttt{balance += amount}, and the compiler generates the
  Cell dictionary operations.
\end{itemize}

%=============================================================
\mysection{Execution Model}\label{sect:exec}
%=============================================================

\mysubsection{Two Execution Paths}

TOS supports two execution paths sharing the same state interface
(\texttt{vm::Dictionary} / \texttt{HashmapE}):

\begin{center}
\begin{tabular}{llll}
\toprule
\textbf{Path} & \textbf{For} & \textbf{Performance} & \textbf{Safety} \\
\midrule
Native C++ & Built-in/system actors & Fastest & Trusted (ships with node) \\
TVM & User-deployed contracts & Medium & TVM sandbox \\
\bottomrule
\end{tabular}
\end{center}

Built-in actors and user-deployed actors can exchange messages freely
because they share the same state format and message protocol.

\mysubsection{TVM Extensions}

Four new TVM opcodes enable actor-native programming:

\begin{center}
\begin{tabular}{ll}
\toprule
\textbf{Opcode} & \textbf{Semantics} \\
\midrule
\texttt{STATEGET} & Read a key from the actor's HashmapE \\
\texttt{STATESET} & Write a key-value pair to the actor's HashmapE \\
\texttt{STATEDEL} & Delete a key from the actor's HashmapE \\
\texttt{ACTORSEND} & Send an asynchronous message to another actor \\
\bottomrule
\end{tabular}
\end{center}

These opcodes are implemented on top of \texttt{vm::Dictionary}
operations---no new data structures are needed. The lower-level
\texttt{DICTGET}/\texttt{DICTSET} instructions also work directly on
the actor's C4 register for advanced use cases.

\mysubsection{Transaction Execution}

The collator dispatches messages to actors via
\texttt{td::actor::send\_closure}. Each \texttt{OnChainActor} is a
\texttt{td::actor::Actor} instance:

\begin{lstlisting}[language=C++,basicstyle=\small\ttfamily]
class OnChainActor : public td::actor::Actor {
    const block::Account* acc_;       // read-only: shared code
    vm::Dictionary actor_state_;      // exclusive: this actor's state
    LogicalTime actor_lt_;            // exclusive: this actor's timeline

    void execute(Ref<vm::Cell> msg,
                 td::Promise<Ref<vm::Cell>> promise) {
        auto result = impl_create_ordinary_transaction(
            msg, acc_, &actor_state_, &actor_lt_, ...);
        promise.set_result(std::move(result));
    }
};
\end{lstlisting}

The collator selects messages, enforces block limits, and assembles
the block. Execution dispatch is asynchronous: the collator sends
actor messages rather than invoking synchronous function calls.

\mysubsection{Execution Phases}

Each actor transaction follows a five-phase execution model:

\begin{enumerate}
\item \textbf{Storage phase} --- compute and deduct storage fees
\item \textbf{Credit phase} --- credit inbound message value
\item \textbf{Compute phase} --- execute contract code (TVM or native)
\item \textbf{Action phase} --- process output actions (messages, reserves)
\item \textbf{Bounce phase} --- generate bounce message on failure
\end{enumerate}

%=============================================================
\mysection{The Tola Language}\label{sect:tola}
%=============================================================

\mysubsection{Design Goals}

Tola is a statically-typed, Solidity-like smart contract language
that compiles to TVM bytecode. Its design goals are:

\begin{enumerate}
\item \textbf{Familiar} --- reuses Solidity syntax conventions so
  developers can onboard in hours, not weeks.
\item \textbf{Actor-native} --- \texttt{actor} replaces
  \texttt{contract}; \texttt{.send()} replaces synchronous
  \texttt{.call()}.
\item \textbf{Safe by default} --- checked arithmetic, no reentrancy
  (impossible by design), no \texttt{delegatecall}, no
  \texttt{tx.origin}.
\item \textbf{100\% FunC coverage} --- every FunC primitive is
  accessible, either natively or through \texttt{unsafe cell \{\}}
  escape blocks.
\end{enumerate}

\mysubsection{Actor Declaration}

\begin{lstlisting}
actor TokenBalance {
    uint64 balance;

    function credit(uint64 amount) external {
        balance += amount;
    }

    function transfer(ActorId recipient, uint64 amount) external {
        require(balance >= amount, "insufficient");
        balance -= amount;
        recipient.send("credit", abi.encode(amount));
    }
}
\end{lstlisting}

Under the hood, \texttt{balance += amount} compiles to
\texttt{STATEGET} $\to$ \texttt{ADD} $\to$ \texttt{STATESET}
operating on the actor's \texttt{HashmapE}. The developer never sees
Cell, Slice, or Builder.

\mysubsection{ActorId and Asynchronous Messaging}

\texttt{ActorId} is a 256-bit first-class type representing an actor's
on-chain identity:

\begin{lstlisting}
ActorId bob = ActorId.derive(self, "bob");
bob.send("credit", abi.encode(100));
// async: fire and forget, result comes back as a separate message
\end{lstlisting}

Send mode flags provide full control over message delivery, covering
all of FunC's \texttt{send\_raw\_message} mode bits:

\begin{lstlisting}
recipient.send("transfer", payload, SendOptions {
    value: 1000,
    mode: SendMode.PAY_FEES_SEPARATELY | SendMode.IGNORE_ERRORS,
    bounce: true
});
\end{lstlisting}

\mysubsection{External Messages, Gas Control, and System Hooks}

Tola covers the full FunC feature set:

\begin{itemize}
\item \texttt{receive\_external} --- handle wallet/dApp messages
  (FunC \texttt{recv\_external})
\item \texttt{gas.accept()} --- accept external messages and pay gas
  (FunC \texttt{accept\_message})
\item \texttt{gas.commit()} --- mid-execution state commit (FunC
  \texttt{commit})
\item \texttt{reserve(amount, mode)} --- balance reservation (FunC
  \texttt{raw\_reserve})
\item \texttt{tick()}/\texttt{tock()} --- system actor lifecycle hooks
  (FunC \texttt{run\_ticktock})
\item \texttt{setCode()} --- on-chain code upgrade (FunC
  \texttt{set\_code})
\item \texttt{crypto.*} --- Ed25519 signatures, SHA-256, keccak-256
\item \texttt{random.*} --- on-chain pseudorandom number generation
\item \texttt{config.param()} --- read blockchain configuration
  parameters
\item \texttt{bounce()} handler --- explicit bounce recovery
\item \texttt{try/catch} --- exception handling (maps to TVM
  \texttt{TRY})
\end{itemize}

\mysubsection{Inheritance and Code Reuse}

Tola supports familiar inheritance patterns via
\texttt{abstract actor} and \texttt{is}:

\begin{lstlisting}
abstract actor Ownable {
    ActorId owner;
    modifier onlyOwner() {
        require(msg.sender == owner, "not owner");
        _;
    }
}

actor TokenBalance is Ownable {
    uint64 balance;
    function transfer(ActorId to, uint64 amount) external onlyOwner {
        require(balance >= amount);
        balance -= amount;
        to.send("credit", abi.encode(amount));
    }
}
\end{lstlisting}

\mysubsection{Raw Cell Escape Hatch}

When interoperating with existing FunC contracts or performing
low-level operations, Tola provides \texttt{unsafe cell} blocks with
full access to Cell/Slice/Builder primitives:

\begin{lstlisting}
function handleLegacy(bytes rawBody) external {
    uint32 op;
    unsafe cell {
        slice s = rawBody.toSlice();
        op = s.loadUint(32);
    }
    if (op == 0x7362d09c) { processTransfer(); }
}
\end{lstlisting}

This ensures 100\% FunC functionality coverage while keeping safe
Tola as the default.

%=============================================================
\mysection{Built-in Standard Actor Library}\label{sect:stdlib}
%=============================================================

\mysubsection{Why Built-in Standard Actors}

A decade of smart contract ecosystem development has demonstrated that
standardized, audited building blocks are the single most impactful
developer tool. The pattern is clear: 90\% of on-chain applications
are compositions of a few standard primitives (tokens, access control,
governance, escrow).

The conventional approach---shipping libraries that each project
copies, modifies, and deploys independently---creates fragmentation,
redundant auditing, and thousands of subtly-different implementations
of the same logic. TOS takes a fundamentally different approach:
standard actors ship as \textbf{native C++ built-ins compiled into the
node binary}:

\begin{center}
\begin{tabular}{lp{5.5cm}p{5.5cm}}
\toprule
& \textbf{Copy-Deploy Model} & \textbf{TOS Built-in Actors} \\
\midrule
Code location & Copied into each project & One copy in node binary \\
Deployment cost & High (per-project bytecode) & Zero (create via message) \\
Audit burden & Each fork audited separately & Audited once, all benefit \\
Fragmentation & Thousands of variants & One canonical implementation \\
Upgrade & Each project independently & Node upgrade, all validators \\
Performance & VM bytecode interpretation & Native C++ execution \\
\bottomrule
\end{tabular}
\end{center}

\mysubsection{Standard Actor Inventory}

Extracted from a decade of production smart contract ecosystem
experience:

\paragraph{Token layer (ships with launch).}
TokenMaster, TokenBalance, TokenAllowance (fungible token standard);
NFTCollection, NFTItem (non-fungible token standard); MultiToken
(multi-asset standard).

\paragraph{Access control layer.}
Ownable, AccessControl (role-based), Multisig (multi-signature),
Timelock (delayed execution).

\paragraph{Governance layer.}
Governor (proposal/voting/execution), VoteActor (voting power
snapshots and delegation).

\paragraph{DeFi layer.}
AMMPool (constant-product swap), LiquidityPosition (concentrated
liquidity), LendingPool (borrow/lend/liquidation), Escrow
(custody/release), Vesting (token vesting schedules), PaymentSplitter
(revenue distribution).

\paragraph{Infrastructure layer.}
Oracle (price feed interface), Proxy (upgradeable actor pattern).

\mysubsection{Usage Model}

Users write only the business logic that differentiates their
application. Standard operations are handled by built-in actors:

\begin{lstlisting}
actor Crowdfund {
    ActorId tokenMaster;
    ActorId treasury;
    uint64 goal;
    uint64 raised;

    function contribute() external payable {
        raised += uint64(msg.value);
        tokenMaster.send("mint_to", abi.encode(
            ActorId.derive(tokenMaster, msg.sender),
            msg.value * 10
        ));
    }

    function finalize() external {
        require(raised >= goal, "goal not met");
        treasury.send("credit", abi.encode(raised));
    }
}
\end{lstlisting}

In this example, the developer writes 15 lines of crowdfunding
logic. Token minting, balance management, and transfer semantics are
handled entirely by the built-in \texttt{TokenMaster} and
\texttt{TokenBalance} actors.

%=============================================================
\mysection{AI Agent Economy}\label{sect:ai}
%=============================================================

The rise of autonomous AI agents creates fundamentally new demands on
blockchain infrastructure. By 2026, AI agents autonomously manage
wallets, negotiate service payments, and execute thousands of
micro-transactions per second. TOS's Actor Model is a natural fit:

\mysubsection{Actors as Agent Wallets}

Each AI agent maps to one or more actors. An agent's spending policy,
authentication credentials, and transaction history are actor state---
isolated, deterministic, and auditable. Unlike account-based chains
where an agent must share state with a global contract, TOS actors
give each agent its own execution context with guaranteed message
ordering.

\begin{lstlisting}
actor AgentWallet is Ownable {
    uint128 dailyLimit;
    uint128 spentToday;
    uint32 dayStart;

    function pay(ActorId service, uint64 amount, bytes memo) external {
        if (block.time > dayStart + 86400) {
            spentToday = 0;
            dayStart = block.time;
        }
        require(spentToday + amount <= dailyLimit, "daily limit");
        spentToday += amount;
        service.send("payment", abi.encode(self, amount, memo),
            SendOptions { value: amount });
    }
}
\end{lstlisting}

\mysubsection{Throughput for Machine-to-Machine Economies}

AI agents produce transaction volumes orders of magnitude higher than
human users. Per-actor parallelism means thousands of agent wallets
execute concurrently---each agent's transactions are independent and
never contend with other agents' state. This is precisely the
throughput model that agent-to-agent payment protocols (x402, MCP)
require.

\mysubsection{Verifiable Agent Behavior}

Because actor state transitions are deterministic and recorded on-chain,
any third party can verify an agent's complete behavioral history.
Combined with ZK proofs (Section~\ref{sect:privacy}), agents can
prove they followed their programmed policy without revealing the
policy itself---critical for competitive trading agents and
confidential service agreements.

%=============================================================
\mysection{Privacy at the Actor Level}\label{sect:privacy}
%=============================================================

\mysubsection{The Privacy Problem in Existing Blockchains}

Current approaches to blockchain privacy fall into two categories:
(1)~privacy-specific chains (Zcash, Monero) that sacrifice programmability,
and (2)~privacy layers bolted onto transparent chains (Aztec on
Ethereum, Tornado Cash) that are awkward, gas-expensive, and
architecturally fragile.

The fundamental issue is that privacy in account-based systems requires
encrypting the \emph{entire} contract state or none of it. There is no
natural boundary for partial confidentiality.

\mysubsection{Actors as Natural Privacy Boundaries}

TOS's Actor Model solves this by making each actor an independent
privacy domain:

\begin{itemize}
\item Each actor's state is a separate HashmapE. Encrypting one
  actor's state entries does not affect any other actor.
\item A \textbf{confidential actor} stores encrypted key-value pairs
  and attaches a zero-knowledge proof to each state transition,
  proving that the transition is valid without revealing the state
  itself.
\item Other actors interact with the confidential actor through the
  same \texttt{.send()} interface---they see the message but not the
  internal state.
\item Privacy is a \emph{per-actor opt-in}, not a network mode.
  Public actors and confidential actors coexist in the same account,
  the same shard, the same block.
\end{itemize}

\mysubsection{ZK-Verified Execution: Verify, Don't Re-Execute}

The actor boundary also enables a shift from \emph{re-execution} to
\emph{verification}:

\begin{enumerate}
\item The actor's owner (or a delegated prover) executes the
  transaction off-chain and produces a succinct ZK proof of correct
  state transition.
\item Validators receive the proof and the new state commitment (the
  HashmapE root hash).
\item Validators \textbf{verify the proof} (a constant-time operation)
  instead of re-executing the computation.
\end{enumerate}

This is architecturally different from ZK rollups, which batch
hundreds of transactions into one proof for a separate L1.  In TOS,
each \emph{individual actor} can independently choose ZK-verified
execution. The result: validators do less work as more actors adopt
ZK proofs, and the network scales by \emph{reducing per-transaction
verification cost} rather than increasing hardware requirements.

A built-in \texttt{ZKVerifier} actor module will provide standard
proof verification for common proof systems (Groth16, PLONK, STARK),
so application developers can attach proofs without implementing
verifiers from scratch.

%=============================================================
\mysection{Post-Quantum Resilience}\label{sect:quantum}
%=============================================================

\mysubsection{The Threat}

Shor's algorithm on a sufficiently powerful quantum computer can break
ECDSA and Ed25519---the signature schemes securing the vast majority
of existing blockchains, including TOS's current default---in
polynomial time. Grover's algorithm weakens symmetric primitives like
SHA-256, though this is addressable by doubling key lengths. The
existential threat is to \emph{public-key cryptography}: any address
whose public key has been exposed on-chain becomes immediately
vulnerable the moment a cryptographically-relevant quantum computer
is operational.

The industry consensus in 2026 is that this threat is 3--10 years
away. NIST has already standardized post-quantum algorithms: CRYSTALS-
Kyber (ML-KEM, FIPS~203) for key encapsulation, CRYSTALS-Dilithium
(ML-DSA, FIPS~204) for digital signatures, and SPHINCS+ (SLH-DSA,
FIPS~205) as a hash-based fallback.

\mysubsection{Actor-Level Cryptographic Agility}

Most blockchains face a dilemma: migrating to post-quantum signatures
requires a network-wide hard fork, new address formats, and mass
migration of user funds---an enormously disruptive operation.

TOS's Actor Model provides a modular alternative:

\begin{itemize}
\item \textbf{Per-actor authentication.} Each actor can independently
  specify its signature scheme. A \texttt{PostQuantumAuth} built-in
  actor module replaces ECDSA/Ed25519 verification with
  SPHINCS+/SLH-DSA for any actor that opts in.
\item \textbf{Gradual migration.} High-value actors (treasuries,
  bridges, governance) can adopt post-quantum signatures immediately.
  Low-value actors can migrate later. No chain-wide hard fork
  required.
\item \textbf{Hybrid signatures.} During the transition period, actors
  can require \emph{both} a classical Ed25519 signature and a
  post-quantum SLH-DSA signature, protecting against both current
  attacks and future quantum attacks simultaneously.
\item \textbf{Algorithm rotation.} If a post-quantum scheme is later
  found to be weak, individual actors can rotate to a newer algorithm
  without affecting the rest of the network.
\end{itemize}

\mysubsection{Hash-Based State Integrity}

TOS's Cell/HashmapE infrastructure is built on hash trees (SHA-256 or
BLAKE2). Against Grover's algorithm, these require only a doubling of
hash output length (256-bit $\to$ 512-bit) to maintain equivalent
security. This upgrade can be performed at the HashmapE level without
changing the actor programming model---developers continue using
\texttt{state\_get}/\texttt{state\_put} regardless of the underlying
hash function.

%=============================================================
\mysection{Concrete Example: Parallel Token Transfers}\label{sect:example}
%=============================================================

Consider a token contract with millions of users. In the account
model, every transfer serializes through one state tree:

\begin{verbatim}
UserA transfer --+
UserB transfer --+--> Token Account (serial)
UserC transfer --+
\end{verbatim}

In TOS's actor model, each user's balance is an independent actor:

\begin{verbatim}
UserA transfer --> BalanceActor(hash(A)) --+
UserB transfer --> BalanceActor(hash(B)) --+--> all parallel
UserC transfer --> BalanceActor(hash(C)) --+
\end{verbatim}

A transfer from Alice to Bob involves two actors:
\begin{enumerate}
\item \texttt{BalanceActor(Alice)} receives the transfer message,
  debits Alice's balance, and sends a \texttt{credit} message to
  \texttt{BalanceActor(Bob)}.
\item \texttt{BalanceActor(Bob)} receives the credit message and
  increments Bob's balance.
\end{enumerate}

Because Alice's and Bob's actors have disjoint state, a simultaneous
transfer from Charlie to Dave executes \textbf{fully in parallel}
with no contention.

%=============================================================
\mysection{Infrastructure Stack}\label{sect:infra}
%=============================================================

TOS's actor-oriented execution model is supported by a full
infrastructure stack designed for high throughput, compact proofs, and
decentralized networking:

\begin{itemize}
\item \textbf{Cell data structure} --- all state and messages are
  encoded as directed acyclic graphs of Cells, providing native Merkle
  proofs at every level of the hierarchy.
\item \textbf{BOC encoding} --- a uniform Bag-of-Cells serialization
  format is used across blocks, network messages, and persistent
  storage.
\item \textbf{Consensus} --- Catchain/Simplex BFT for fast finality
  with Byzantine fault tolerance.
\item \textbf{Networking} --- ADNL (Abstract Datagram Network Layer),
  RLDP (Reliable Large Datagram Protocol), a Kademlia-based DHT, and
  QUIC transport.
\item \textbf{Dynamic sharding} --- address-prefix-based shard
  splitting allows the network to scale horizontally as load increases.
\item \textbf{Storage} --- CellDB backed by RocksDB for persistent
  state.
\item \textbf{Execution phases} --- storage $\to$ credit $\to$
  compute $\to$ action $\to$ bounce, applied per actor transaction.
\item \textbf{FunC/Tolk compatibility} --- contracts written in FunC
  or Tolk run on TVM alongside Tola actors, sharing the same message
  protocol.
\item \textbf{Lite-client / Tonlib} --- lightweight client tools for
  querying chain state and submitting transactions.
\end{itemize}

The actor model determines the \emph{granularity} of execution (the
actor, not the account, is the scheduling unit) and the
\emph{developer interface} (key-value state rather than raw
Cell/Slice/Builder manipulation). The infrastructure layers below
this boundary are optimized for the actor workload from the start.

%=============================================================
\mysection{Comparison with Existing Platforms}\label{sect:compare}
%=============================================================

\begin{center}
\small
\begin{tabular}{lccc}
\toprule
& \textbf{TOS} & \textbf{Serial-EVM Chains} & \textbf{Object/Resource Chains} \\
\midrule
Parallel execution & Actor-level & Planned / L2 only & Object/resource \\
Developer language & Tola (Solidity-like) & Solidity & New language (Move) \\
Learning curve & Low & Baseline & High \\
Reentrancy & Impossible & Major risk & N/A \\
Standard library & Native built-in & Copy-deploy per project & None standard \\
Hot-contract scaling & $M\times$ (per-actor) & $1\times$ & $M\times$ (per-object) \\
Privacy model & Per-actor ZK opt-in & L2 rollups only & None \\
Post-quantum & Per-actor migration & Global hard fork & Global hard fork \\
AI agent support & Native actor wallets & Add-on abstraction & None standard \\
ZK verification & Per-actor verify & L2 proof batching & None \\
Merkle proofs & Cell (built-in) & Trie-based & None \\
\bottomrule
\end{tabular}
\end{center}

Existing platforms address parallelism, privacy, quantum resilience,
and agent support as separate problems with separate, layered
solutions---rollups for scaling, privacy protocols for confidentiality,
hard forks for cryptographic migration, account abstraction for agents.
TOS starts from first principles: the actor is the unit of execution,
privacy, and cryptographic identity. Parallelism, confidentiality,
and quantum resilience are not features added to the architecture---
they are \emph{consequences} of it.

%=============================================================
\mysection{Roadmap}\label{sect:roadmap}
%=============================================================

\begin{enumerate}
\item \textbf{Phase 1: Actor State Isolation + Built-in Actors.}
  Add actor sub-structure to accounts, implement TVM opcodes
  (\texttt{STATEGET}/\texttt{STATESET}/\texttt{ACTORSEND}), build
  native Token/NFT actors, validate architecture end-to-end.

\item \textbf{Phase 2: Per-Actor Parallel Execution.}
  Each on-chain actor becomes a \texttt{td::actor::Actor}. Collator
  dispatches via \texttt{send\_closure}. Different actors under the
  same account execute in parallel.

\item \textbf{Phase 3: Message Routing + Tola Compiler.}
  Add \texttt{actor\_id} to \texttt{OutMsgQueue} key. Ship the Tola
  compiler (Solidity-like syntax $\to$ TVM bytecode). Full standard
  actor library (DeFi, governance, access control).

\item \textbf{Phase 4: Privacy Actors + ZK Verification.}
  Built-in \texttt{ConfidentialActor} and \texttt{ZKVerifier} modules.
  Per-actor opt-in encrypted state with ZK proof of correct state
  transition. Validators verify proofs instead of re-executing.

\item \textbf{Phase 5: Post-Quantum Authentication.}
  Built-in \texttt{PostQuantumAuth} module (SPHINCS+/SLH-DSA).
  Per-actor signature scheme migration. Hybrid classical+post-quantum
  mode for the transition period.

\item \textbf{Phase 6: AI Agent Economy.}
  Built-in \texttt{AgentWallet} with spending policies, daily limits,
  allowlists. Native integration with x402/MCP payment protocols.
  High-throughput agent-to-agent settlement.
\end{enumerate}

%=============================================================
\section*{Conclusion}
\markbothsame{\textsc{Conclusion}}
\addcontentsline{toc}{section}{Conclusion}
%=============================================================

Blockchains in 2026 face simultaneous demands for higher throughput,
AI agent support, user privacy, and quantum resilience. Existing
platforms address these as separate problems with separate solutions---
rollups for scaling, privacy layers for confidentiality, hard forks
for post-quantum migration, account abstraction for agents. The
result is architectural complexity that grows with each new
requirement.

TOS takes a different approach. The Actor Model is not merely a
scaling technique---it is an \emph{organizing principle} from which
parallelism, privacy, quantum resilience, and agent support all emerge
naturally:

\begin{itemize}
\item \textbf{Parallelism.} Hot contracts scale linearly with the number
  of independent actors---a direct consequence of state isolation.
\item \textbf{Privacy.} Each actor is an independent privacy domain.
  Confidential actors encrypt their state and prove transitions via ZK
  proofs---a direct consequence of actor boundaries.
\item \textbf{Quantum resilience.} Each actor independently adopts
  post-quantum signatures---a direct consequence of actor-level
  authentication.
\item \textbf{AI agents.} Each agent is an actor with its own wallet,
  policy, and message interface---a direct consequence of the actor as
  the minimum execution unit.
\item \textbf{Verification, not re-execution.} Validators verify ZK
  proofs instead of re-computing state transitions---a direct
  consequence of per-actor proof attachment.
\item \textbf{Developer simplicity.} Solidity-like syntax, one new
  concept (\texttt{.send()} instead of \texttt{.call()}), and 90\% of
  common functionality built-in and free to use.
\item \textbf{Production-grade foundation.} Consensus, networking,
  storage, and Merkle proofs built on proven, well-understood
  techniques---Cell trees, BFT consensus, Kademlia DHT---chosen
  specifically to complement the actor execution model.
\end{itemize}

The key architectural insight is that Cell infrastructure need not be
replaced---only \emph{encapsulated}---and the actor boundary is not
just a concurrency boundary but also a privacy boundary, a
cryptographic boundary, and an economic boundary. One abstraction,
many consequences.

%=============================================================
\clearpage
\markbothsame{\textsc{References}}
\begin{thebibliography}{99}

\bibitem{Agha86}
  {\sc G.~Agha}, {\sl Actors: A Model of Concurrent Computation in Distributed Systems},
  MIT Press, 1986.

\bibitem{Hewitt73}
  {\sc C.~Hewitt, P.~Bishop, R.~Steiger}, {\sl A Universal Modular ACTOR Formalism for Artificial Intelligence},
  IJCAI, 1973.

\bibitem{TONOriginal}
  {\sc N.~Durov}, {\sl Telegram Open Network},
  2018.

\bibitem{Wood14}
  {\sc G.~Wood}, {\sl A Secure Decentralised Generalised Transaction Ledger},
  2014 (updated 2024). Foundational work on serial-execution smart contract VMs.

\bibitem{MoveLanguage}
  {\sc S.~Blackshear et al.}, {\sl Move: A Language With Programmable Resources},
  2020. Resource-oriented approach to parallel smart contract execution.

\bibitem{PBFT}
  {\sc M.~Castro, B.~Liskov}, {\sl Practical Byzantine Fault Tolerance},
  OSDI, 1999.

\bibitem{Kademlia}
  {\sc P.~Maymounkov, D.~Mazi\`{e}res}, {\sl Kademlia: A Peer-to-peer Information System Based on the XOR Metric},
  IPTPS, 2002.

\bibitem{x402}
  {\sc Linux Foundation}, {\sl x402: HTTP Payment Protocol for Agentic Commerce},
  2026. Adopted by Visa, Stripe, AWS, Anthropic, and Google Cloud for
  machine-to-machine payments.

\bibitem{MCP}
  {\sc Anthropic}, {\sl Model Context Protocol (MCP): Standardizing Agent Tool Invocation},
  2025.

\bibitem{QuantumThreat}
  {\sc Google Quantum AI}, {\sl How a Quantum Computer Could Break Blockchain Cryptography},
  reported in CoinDesk, April 2026.

\bibitem{SPHINCS}
  {\sc D.~Bernstein et al.}, {\sl SPHINCS+: Stateless Hash-Based Signatures},
  NIST FIPS~205 (SLH-DSA), standardized August 2024.

\bibitem{Dilithium}
  {\sc L.~Ducas et al.}, {\sl CRYSTALS-Dilithium: Digital Signatures from Module Lattices},
  NIST FIPS~204 (ML-DSA), standardized August 2024.

\bibitem{Groth16}
  {\sc J.~Groth}, {\sl On the Size of Pairing-Based Non-interactive Arguments},
  EUROCRYPT, 2016.

\bibitem{PLONK}
  {\sc A.~Gabizon, Z.~Williamson, O.~Ciobotaru}, {\sl PLONK: Permutations over Lagrange-Bases for Oecumenical Noninteractive Arguments of Knowledge},
  IACR ePrint 2019/953.

\bibitem{ZKRollup}
  {\sl Zero-Knowledge Rollups for Private Smart Contract Execution},
  various implementations, 2023--2026.

\bibitem{AcctAbstraction}
  {\sl Account Abstraction for Programmable Wallets},
  industry standard, 2023--2026.

\end{thebibliography}

\end{document}
